\documentclass[11pt,a4paper]{article}

\usepackage[margin=2.4cm]{geometry}
\usepackage{amsmath,amssymb,amsthm}
\usepackage{booktabs}
\usepackage{microtype}
\usepackage{enumitem}
\usepackage[hidelinks]{hyperref}

\newtheorem{definition}{Definition}[section]
\newtheorem{proposition}[definition]{Proposition}
\newtheorem{lemma}[definition]{Lemma}
\theoremstyle{remark}
\newtheorem{remark}[definition]{Remark}

\newcommand{\obs}{\mathrm{obs}}
\newcommand{\adj}{\bar p}
\newcommand{\lift}{\operatorname{lift}}
\newcommand{\rep}{\operatorname{rep}}

\title{How PaM-ST computes and corrects its $p$-values}
\author{}
\date{22 September 2026}

\begin{document}
\maketitle

\begin{abstract}
A self-contained account of the inference step: the table of counts over
candidates and null model instances, the rank-based raw $p$-value read along
its rows, the family minimum read down its columns, and the Westfall--Young
step-down correction built from the two. Every number in the worked examples
comes from the four-tissue run in \texttt{output/minp\_all}
($C=2{,}013$ candidates, $B=40{,}000$ instances). Companion document:
\texttt{docs/pipeline.tex}, which covers the whole pipeline; this note expands
its Sections on the $p$-value and the multiplicity correction.
\end{abstract}

% ---------------------------------------------------------------
\section{The table of counts}

\begin{definition}[Candidates, instances, counts]
\label{def:table}
Let $c=1,\dots,C$ index the candidate patterns, fixed on the selection half of
the checkerboard before any shuffling, and let $b=1,\dots,B$ index the null
model instances: independent labellings drawn by permuting cell-type labels
within each block tile of the testing half. Write $b=0$ for the observed
tissue. Let
\[
  N_c^{(b)}\in\mathbb{Z}_{\ge0}
  \quad\text{be the number of testing circles matching candidate $c$ under labelling $b$,}
\]
so the inference operates on a $C\times(B+1)$ table of non-negative integers,
whose column $b=0$ is the data and whose remaining $B$ columns are chance.
Write $N_c^{\obs}=N_c^{(0)}$.
\end{definition}

The two readings of this table give the two quantities of interest: rows give
raw $p$-values (Section~\ref{sec:raw}), columns give the null distribution of
the best candidate in the family (Section~\ref{sec:adj}).

% ---------------------------------------------------------------
\section{Rows: the raw \texorpdfstring{$p$}{p}-value}
\label{sec:raw}

\begin{definition}[Rank and score]
\label{def:score}
For candidate $c$ define the upper-tail rank within its own row,
\[
  R_c(y)=\bigl|\{b\in\{0,1,\dots,B\}: N_c^{(b)}\ge y\}\bigr|,
\]
counting the observed column and repetitions. The \emph{score} of column $b$
for candidate $c$ is
\[
  s_c^{(b)}=\frac{R_c\bigl(N_c^{(b)}\bigr)}{B+1}\in\Bigl[\tfrac1{B+1},\,1\Bigr].
\]
The \emph{raw $p$-value} of candidate $c$ is the score of the observed column,
\[
  p_c:=s_c^{(0)}=\frac{R_c(N_c^{\obs})}{B+1}
     =\frac{1+\bigl|\{b\ge1: N_c^{(b)}\ge N_c^{\obs}\}\bigr|}{B+1}.
\]
\end{definition}

Three properties follow immediately and matter in practice.

\begin{enumerate}[itemsep=3pt]
  \item \textbf{One-sided.} Only $\ge$ is counted: a candidate that occurs
    \emph{less} often than chance receives $p_c=1$, not a small $p$-value.
  \item \textbf{Floor.} $p_c\ge 1/(B+1)$, here $1/40{,}001=2.5\cdot10^{-5}$;
    "never matched in $B$ instances" cannot produce $p_c=0$.
  \item \textbf{Ties are conservative.} Counts are integers, and tied values
    all receive the larger rank, which can only increase $p_c$.
\end{enumerate}

\begin{proposition}[Validity]
\label{prop:valid}
If $N_c^{(0)},\dots,N_c^{(B)}$ are exchangeable --- which holds under the
conditional block null, since the observed labelling is then one more uniform
draw from the same set of rearrangements --- then
$\Pr(p_c\le\alpha)\le\alpha$ for every $\alpha\in(0,1)$.
\end{proposition}

\begin{proof}
Exchangeability makes the rank of $N_c^{(0)}$ uniform on $\{1,\dots,B+1\}$,
ties only increasing $R_c$. Hence $\Pr(R_c(N_c^{(0)})\le k)\le k/(B+1)$ for
every integer $k$; take $k=\lfloor\alpha(B+1)\rfloor$.
\end{proof}

\begin{remark}[Why scores, not counts]
\label{rem:scale}
The scores put every candidate on one scale, which is what makes the minimum in
Section~\ref{sec:adj} meaningful. Raw counts do not: in this run the reported
candidates have $N_c^{\obs}$ between $73$ and $1{,}497$ and null means between
$0.01$ and $383$, so a count of $200$ is overwhelming evidence for one
candidate and below average for another. The map $N\mapsto s$ is exactly the
within-row calibration that removes those scale differences.
\end{remark}

\begin{remark}[The observed column belongs in the pool]
\label{rem:pool}
$R_c$ counts over $\{0,1,\dots,B\}$ for the scores of instances too, so an
instance is ranked against the same $B+1$ values as the observation. Scoring
instances against the other instances only would give them denominator $B$ and
floor $1/B$, while the observation reaches $1/(B+1)$: an observation at its
floor could then never be matched, and in simulation the family-wise error rate
rose from the nominal $5\%$ to $10$--$16\%$.
\end{remark}

% ---------------------------------------------------------------
\section{Columns: the family minimum and the correction}
\label{sec:adj}

\begin{definition}[Family minimum]
For instance $b\ge1$,
\[
  M^{(b)}=\min_{1\le c\le C}s_c^{(b)}
\]
is the score of the most extreme candidate anywhere in the family on that
instance: what the search would have reported had that instance been the data.
\end{definition}

\begin{definition}[Single-step adjusted $p$-value]
\label{def:single}
\[
  \adj^{\,\mathrm{ss}}_c=\frac{1+\bigl|\{b\ge1: M^{(b)}\le p_c\}\bigr|}{B+1}.
\]
\end{definition}

The raw $p$-value enters here as the \emph{threshold}: the correction asks how
often chance, anywhere in the family, produced a result at least as extreme as
this candidate's own.

\begin{definition}[Step-down with monotone enforcement]
\label{def:stepdown}
Order the candidates by increasing raw $p$-value, ties broken by decreasing
lift and then by index, and write $(1),\dots,(C)$ for that order, so
$p_{(1)}\le\dots\le p_{(C)}$. For $k=1,\dots,C$ put
\[
  M^{(b)}_k=\min_{j\ge k}s_{(j)}^{(b)},
  \qquad
  \tilde p_{(k)}=\frac{1+\bigl|\{b\ge1:M^{(b)}_k\le p_{(k)}\}\bigr|}{B+1},
  \qquad
  \adj_{(k)}=\max_{j\le k}\tilde p_{(j)}.
\]
Candidate $c$ is declared significant when $\adj_c\le\alpha$.
\end{definition}

\begin{proposition}[Family-wise error]
Under the conditional null and subset pivotality --- the joint law of
$(p_c)_{c\in I_0}$, with $I_0$ the set of true nulls, is the same as it would
be if every hypothesis were true --- the procedure of
Definition~\ref{def:stepdown} satisfies
$\Pr(\exists\,c\in I_0:\adj_c\le\alpha)\le\alpha$.
\end{proposition}

This is the Westfall--Young result: the instances supply the joint distribution
of the scores, dependence between overlapping candidates included, so nothing
is assumed about independence.

\begin{proposition}[Step-down is never worse]
$\adj_c\le\adj^{\,\mathrm{ss}}_c$ for every $c$.
\end{proposition}

\begin{proof}
$M^{(b)}_k$ is a minimum over a subset of the candidates minimised in
$M^{(b)}$, so $M^{(b)}_k\ge M^{(b)}$ and hence $\tilde p_{(k)}\le
\adj^{\,\mathrm{ss}}_{(k)}$. The running maximum preserves this because
$\adj^{\,\mathrm{ss}}_{(j)}$ is non-decreasing in $j$: $p_{(j)}$ is
non-decreasing and Definition~\ref{def:single} is monotone in its threshold.
\end{proof}

\begin{lemma}[Ties at the floor share one adjusted value]
\label{lem:ties}
If $p_{(1)}=\dots=p_{(m)}$, then $\adj_{(1)}=\dots=\adj_{(m)}$.
\end{lemma}

\begin{proof}
For $k\le m$ the threshold $p_{(k)}$ is constant while $M_k^{(b)}$ is
non-decreasing in $k$, so $\tilde p_{(k)}$ is non-increasing in $k$; the
running maximum therefore holds every $\adj_{(k)}$ at $\tilde p_{(1)}$.
\end{proof}

Lemma~\ref{lem:ties} explains a feature of the output that is otherwise
puzzling: $38$ of the reported patterns carry exactly $\adj=0.00715$. They all
sit at the floor $p_c=1/(B+1)$, so no ordering among them survives the
correction, and they can only be ranked by effect size.

% ---------------------------------------------------------------
\section{The four uses of the raw \texorpdfstring{$p$}{p}-value}

\begin{enumerate}[itemsep=3pt]
  \item \textbf{Threshold} in Definitions~\ref{def:single}
    and~\ref{def:stepdown}: the correction counts instances whose family
    minimum is at most $p_c$.
  \item \textbf{Ordering}: the step-down sequence, and therefore which
    restricted family each candidate faces, is the ordering by $p_c$.
  \item \textbf{Per-sample evidence}: with $S$ samples, the same construction
    restricted to sample $s$ gives $p_{c,s}$, and the reported replication
    count is $\rep_c=|\{s:p_{c,s}\le\alpha\}|$. These are raw values, not
    adjusted, and carry no family-wise guarantee.
  \item \textbf{Diagnostics}: the mass of raw $p$-values at the floor measures
    how far the permutation budget limits resolution, and
    $\min_c\adj_c\gtrsim C/(B+1)$ is the rule behind $B\gtrsim C/\alpha$.
\end{enumerate}

The raw $p$-value is never the decision. In the run below, $51$ of the $60$
reported patterns satisfy $p_c\le0.05$ while $41$ satisfy $\adj_c\le0.05$.

% ---------------------------------------------------------------
\section{Worked example from the four-tissue run}

$C=2{,}013$, $B=40{,}000$, $\alpha=0.05$. Five rows of the table, with the
observed column and the first six instances:

\begin{center}
\small
\begin{tabular}{lrrrrrrrrr}
\toprule
Candidate & $N^{\obs}$ & $b{=}1$ & $2$ & $3$ & $4$ & $5$ & $6$ & $p_c$ & $\adj_c$ \\
\midrule
Cancer nest $32{:}3$        & 221 & 1 & 2 & 1 & 0 & 3 & 0 & $1/40{,}001$ & $0.0071$ \\
Vessels $22{:}9$            & 122 & 24 & 22 & 25 & 20 & 24 & 32 & $1/40{,}001$ & $0.0071$ \\
Stroma $+$ mem.\ CD4 $15{:}2$ & 266 & 198 & 157 & 204 & 180 & 170 & 142 & $7/40{,}001$ & $0.0360$ \\
Stroma $+$ vessels $23{:}2$ & 304 & 234 & 196 & 221 & 189 & 241 & 189 & $13/40{,}001$ & $0.0606$ \\
Cancer $+$ stroma $11{:}10$ & 182 & 346 & 327 & 335 & 318 & 324 & 326 & $1$ & $1$ \\
\bottomrule
\end{tabular}
\end{center}

Reading the rows. For the vessel pattern the $40{,}000$ instance counts have
mean $28.9$ and standard deviation $8.0$, a maximum of $73$, and none reaches
$N^{\obs}=122$; hence $R_c(122)=1$ and $p_c=1/40{,}001$. Its effect sizes are
$\lift_c=122/28.9=4.2$ and $z_c=(122-28.9)/8.0=11.6$. For the last row the
observed value sits below every instance, so $R_c(182)=40{,}001$ and $p_c=1$:
the pattern is rarer in the tissue than chance would make it.

Reading the columns. On $285$ of the $40{,}000$ instances some candidate
reached the floor, i.e.\ $M^{(b)}_1\le 1/40{,}001$, so
\[
  \adj_{(1)}=\frac{1+285}{40{,}001}=\frac{286}{40{,}001}=0.00715 .
\]
Further down the ranking the thresholds are looser and the families smaller:
\[
  \adj_{(k)}=\frac{1{,}442}{40{,}001}=0.0360
  \quad\text{at } p_{(k)}=\frac{7}{40{,}001},
  \qquad
  \adj_{(k)}=\frac{2{,}424}{40{,}001}=0.0606
  \quad\text{at } p_{(k)}=\frac{13}{40{,}001}.
\]
The first is significant, the second is not, and the boundary of the reported
list falls between them: $41$ of the $60$ reported patterns satisfy
$\adj_c\le0.05$, of which $38$ are the tied block of Lemma~\ref{lem:ties}.

% ---------------------------------------------------------------
\section{Budget}

Since every score is a multiple of $1/(B+1)$, so is every adjusted value, and
if the candidates behaved independently a fraction $\approx C/(B+1)$ of the
instances would contain some candidate at its floor. Hence
\[
  \min_c\adj_c\ \gtrsim\ \frac{C}{B+1},
  \qquad\text{so rejecting at level }\alpha\text{ needs } B\gtrsim C/\alpha .
\]
Here $C/(B+1)=2{,}013/40{,}001=0.050$, while the observed minimum is
$0.00715$, about seven times smaller: the candidates overlap heavily, so the
family behaves as if it were far smaller than $2{,}013$ independent tests.

\end{document}
